\section{Mathematical construction on curvilinear elements}
\label{sec:curvilinear-mathematics}
The affine implementation relies on two simplifications: the Jacobian is
constant within an element, and differentiation stays inside the represented
polynomial space. A curved map removes both simplifications. Our extension
keeps the physical projection idea of OFDG, but assembles its operators for
each curved element and projects each derivative back into the solution space.
This section extends the coefficient construction in
\cref{sec:ofdg-projection-coefficients,sec:ofdg-coefficient-derivatives,sec:ofdg-face-evaluation}
to nonconstant geometric transformations.

\subsection{What is represented, and what is differentiated?}
The input to the derivative calculation is the DG function on one element,
not the unknown exact solution of the conservation law. Its coefficient
vector $U=(U_1,\ldots,U_{N_k})^{\mathsf T}$ represents
\begin{equation}
 u_h(x)=\sum_{b=1}^{N_k}U_b\phi_b(x),\qquad
 \widehat u_h(\xi)=u_h(T_K(\xi))
                 =\sum_{b=1}^{N_k}U_b\widehat\phi_b(\xi).
 \label{eq:curved-stored-function}
\end{equation}
The coefficients are the same in these two expressions. The physical basis
is obtained by composing the reference basis with $T_K^{-1}$, so geometry
changes the function's dependence on physical position without changing the
meaning of the coefficient expansion. In particular, a polynomial in $\xi$
need not be a polynomial in $x$.

All differentiations below are taken inside one element. We do not take a
distributional derivative across the discontinuities of the DG solution.
The separate element derivatives are evaluated at faces and their jumps are
formed afterwards, as in \cref{sec:ofdg-face-evaluation}. This distinction
allows each element to have a well-defined smooth interior derivative even
when the global DG function is discontinuous.

Reference differentiation is already available through
$D_{\xi_r}=V^{-1}G^{(r)}$ from
\cref{sec:ofdg-coefficient-derivatives}. The next question is how to turn
this reference operation into a physical derivative when the element map is
non-affine. We first identify the representation problem, then construct
one projected derivative, and finally explain repeated differentiation.

\subsection{Why the affine construction is insufficient}
Let $\boldsymbol x=T_K(\boldsymbol\xi)$ be a smooth, invertible,
orientation-preserving element map, and write
\begin{equation}
 J_K(\boldsymbol\xi)=\frac{\partial T_K}{\partial\boldsymbol\xi},
 \qquad w_K(\boldsymbol\xi)=\det J_K(\boldsymbol\xi)>0.
 \label{eq:curved-jacobian}
\end{equation}
For $v=\widehat v\circ T_K^{-1}$, change of variables and the chain rule give
\begin{align}
 \int_K v\,dx&=\int_{\widehat K}\widehat v\,w_K\,d\xi,
 \label{eq:curved-change-variables}\\
 (\nabla_xv)\circ T_K&=J_K^{-\mathsf T}\nabla_\xi\widehat v.
 \label{eq:curved-chain-rule}
\end{align}
The physical inner product therefore has a nonconstant weight. Even if
$\widehat v$ and $T_K$ are polynomial, the inverse Jacobian generally makes
$\partial_{x_j}v$ fall outside the mapped finite element space. Evaluating a
single inverse Jacobian at the element centre cannot represent this variation.
Componentwise, the chain rule reads
\begin{equation}
 (\partial_{x_j}u_h)(T_K(\xi))
 =\sum_{r=1}^{d}(J_K^{-1}(\xi))_{rj}
                  \sum_{b=1}^{N_k}U_b\partial_{\xi_r}\widehat\phi_b(\xi).
\end{equation}
Thus the exact first physical derivative of the represented $u_h$ can still
be evaluated at any quadrature point using the available basis derivatives
and Jacobian. The difficulty is storing that entire derivative with only
$N_k$ coefficients in the original space. On an affine map the geometric
factors are constant and the representation is exact; on a curved map they
are position-dependent and generally take the derivative outside the space.

A one-dimensional example makes the obstruction explicit. On $[0,1]$, take
\begin{equation}
 x=T(\xi)=\xi+\varepsilon\xi(1-\xi),\quad
 0<|\varepsilon|<1,\qquad \widehat v(\xi)=\xi.
\end{equation}
Although the reference state is linear,
\begin{equation}
 (\partial_xv)\circ T=\frac{1}{1+\varepsilon-2\varepsilon\xi},\qquad
 (\partial_x^2v)\circ T=
 \frac{2\varepsilon}{(1+\varepsilon-2\varepsilon\xi)^3}.
 \label{eq:curved-derivative-example}
\end{equation}
Neither expression is generally a finite-degree reference polynomial. The
issue is already present for an ordinary polynomial geometry map; it does not
require NURBS.

\Cref{fig:curved-map} shows the corresponding geometric issue in two dimensions:
reference coordinate lines are mapped to a curved physical element, and the
Jacobian weight and face normals vary within that element.
\begin{figure}[htbp]
\centering
\begin{tikzpicture}[>=Latex,font=\small]
 \node at (1.4,4.75) {reference element $\widehat K$};
 \node at (7.4,4.75) {curved physical element $K$};
 \foreach \s in {0,.25,.5,.75,1} {
  \draw[blue!30] ({2.8*\s},0)--({2.8*\s},2.8);
  \draw[blue!30] (0,{2.8*\s})--(2.8,{2.8*\s});
  \draw[blue!30,domain=0:1,samples=35]
    plot ({6+2.8*\x},{2.8*((1+.25*\x)*\s+.2*\x*\x)});
  \draw[blue!30] ({6+2.8*\s},{2.8*.2*\s*\s})--
    ({6+2.8*\s},{2.8*(1+.25*\s+.2*\s*\s)});
 }
 \draw[blue!70!black,thick] (0,0) rectangle (2.8,2.8);
 \foreach \s in {0,1} {
  \draw[blue!70!black,thick,domain=0:1,samples=35]
    plot ({6+2.8*\x},{2.8*((1+.25*\x)*\s+.2*\x*\x)});
  \draw[blue!70!black,thick] ({6+2.8*\s},{2.8*.2*\s*\s})--
    ({6+2.8*\s},{2.8*(1+.25*\s+.2*\s*\s)});
 }
 \foreach \s in {.2,.5,.8} {
  \draw[->,orange!85!black,thick]
    ({6+2.8*\s},{2.8*(1+.25*\s+.2*\s*\s)})--
    ({6+2.8*\s-.52*(.25+.4*\s)},{2.8*(1+.25*\s+.2*\s*\s)+.52});
 }
 \draw[->,thick] (3.25,1.55)--(5.4,1.55);
 \node[above] at (4.3,1.55) {$T_K$};
 \node[below,align=center] at (1.4,-.25) {straight reference grid\\coordinates $(\xi,\eta)$};
 \node[below,align=center] at (7.4,-.25) {mapped grid\\spatially varying Jacobian and normals};
\end{tikzpicture}
\caption{An illustrative polynomial map
$T_K(\xi,\eta)=(\xi,(1+\xi/4)\eta+\xi^2/5)$ from the unit square.
Its determinant is $1+\xi/4>0$: equal reference areas need not have equal
physical areas. Orange arrows show outward normal directions along the curved
top face (not their magnitudes). The interior lines are coordinate lines
within one element, not separate DG cells. This is a geometric schematic,
not a mesh from a numerical experiment.}
\label{fig:curved-map}
\end{figure}


\subsection{Physical mass matrices and nested projections}
Use the mapped spaces $V^r(K)$ from \cref{eq:ofdg-space}, all with the same
map $T_K$. Let $\widehat\phi_i$ span the order-$k$ reference space and
$\widehat\psi_a^{(r)}$ its order-$r$ subspace. Their physical mass and coupling
matrices are defined below. Let $N_k$ and $N_r$ be the numbers of basis
functions at orders $k$ and $r$. Then $M_k^K$ is $N_k\times N_k$,
$M_r^K$ is $N_r\times N_r$, and the coupling matrix $B_r^K$ is
$N_r\times N_k$. Indices $i,j$ run over the order-$k$ basis and $a,b$
over the order-$r$ basis:
\begin{align}
 (M_k^K)_{ij}&=\int_{\widehat K}\widehat\phi_i\widehat\phi_j w_K\,d\xi,\\
 (M_r^K)_{ab}&=\int_{\widehat K}
       \widehat\psi_a^{(r)}\widehat\psi_b^{(r)}w_K\,d\xi,\\
 (B_r^K)_{ai}&=\int_{\widehat K}
       \widehat\psi_a^{(r)}\widehat\phi_iw_K\,d\xi.
 \label{eq:curved-weighted-masses}
\end{align}
If $U$ represents $u_h$, its physical $L^2$ projection into $V^r(K)$ has
lower-order coefficients $C_r$ determined by
\begin{equation}
 M_r^K C_r=B_r^K U.
 \label{eq:curved-projection-solve}
\end{equation}
This is the usual finite element projection equation, now with the physical
Jacobian weight retained. To express the projected function in the original
order-$k$ basis, solve $M_k^K U_r=(B_r^K)^{\mathsf T}C_r$. Thus the matrix
representing $\Pi_K^r$ in that basis is
\begin{equation}
 P_r^K=(M_k^K)^{-1}(B_r^K)^{\mathsf T}(M_r^K)^{-1}B_r^K.
 \label{eq:curved-projector-matrix}
\end{equation}
The inverse notation denotes the associated linear solves. Nesting makes
this lifting exact as a function-space operation. With exact integrals,
\begin{equation}
 P_r^KP_s^K=P_{\min(r,s)}^K,\qquad
 (P_r^K)^{\mathsf T}M_k^K=M_k^KP_r^K.
 \label{eq:curved-projector-identities}
\end{equation}
Consequently the physical orthogonal shells and the frozen exponential
attenuation in \cref{eq:ofdg-exact-decay} remain available. No orthogonal
storage basis is required.

Mean preservation follows directly from testing the projection equation with
one: $\int_K(\Pi_K^r u_h-u_h)\,dx=0$. In coefficient form, let
$b_i=\int_K\phi_i\,dx$; then $b^{\mathsf T}P_r^K=b^{\mathsf T}$. Every
nonconstant shell has zero $b$-weighted integral, so damping it leaves
$b^{\mathsf T}U$ unchanged component by component. On affine elements,
$w_K$ is constant and cancels from \cref{eq:curved-projector-matrix}, recovering
the reusable reference projector. On curved elements it cannot be cancelled.

\subsection{Projecting a physical derivative}
Let $f=\partial_{x_j}u_h$ denote the exact physical derivative of the
represented input. Since $f$ need not belong to $V^k(K)$, we select the
representable function with the smallest physical squared-integral error:
\begin{equation}
 z_h=\underset{z\in V^k(K)}{\operatorname{argmin}}\,
           \int_K |z-f|^2\,dx
     =\Pi_K^k f=\mathscr D_{K,j}u_h.
 \label{eq:curved-derivative-minimization}
\end{equation}
The word physical refers to the measure $dx$. A region stretched by the
map occupies more physical volume and must have correspondingly more weight
in this error. Minimizing an unweighted reference integral would generally
choose a different approximation.

To obtain the minimizer, require the first variation in every available test
direction $v_h$ to vanish. This gives the orthogonality condition
\begin{equation}
 (z_h,v_h)_K=(\partial_{x_j}u_h,v_h)_K
 \quad\text{for every }v_h\in V^k(K).
 \label{eq:curved-derivative-weak}
\end{equation}
The residual $z_h-f$ is therefore orthogonal to the available approximation
space. This does not require it to vanish at nodes or face points: physical
$L^2$ projection is an integral best approximation, not interpolation.
Write the unknown projected derivative as
$z_h=\sum_b Z_b\phi_b$ and test with each $\phi_a$. The result is
$\sum_b(\phi_a,\phi_b)_K Z_b=
 \sum_b(\phi_a,\partial_{x_j}\phi_b)_K U_b$.
The left-hand side is the physical mass matrix acting on $Z$; the
right-hand side supplies the inner products with the known derivative.
Using \cref{eq:curved-chain-rule} gives
\begin{align}
 M_k^K Z&=G_{K,j}U,\qquad D_{K,j}=(M_k^K)^{-1}G_{K,j},\\
 (G_{K,j})_{ab}&=\int_{\widehat K}\widehat\phi_a
  \left[\sum_{s=1}^d(J_K^{-1})_{sj}
                   \partial_{\xi_s}\widehat\phi_b\right]w_K\,d\xi.
 \label{eq:curved-derivative-assembly}
\end{align}
Here $U,Z\in\mathbb R^{N_k}$ and $M_k^K$, $G_{K,j}$, and $D_{K,j}$
are all $N_k\times N_k$. The subscript $j$ selects one physical coordinate
direction. The integral matrix $G_{K,j}$ must not be confused with the
sampled reference-derivative matrix $G^{(r)}$ used to construct $D_{\xi_r}$.
The matrix $D_{K,j}$ maps input coefficients directly to projected derivative
coefficients: $Z=D_{K,j}U$.

At each volume quadrature point, assembly evaluates the reference basis,
its first derivatives, the inverse Jacobian, and the physical weight. It
accumulates the mass and derivative inner products, then solves the mass
system for each column of $G_{K,j}$. The static geometry allows the resulting
$D_{K,j}$ to be stored and reused. During filtering, applying a derivative is
a matrix multiplication; no new global PDE problem is solved.

The geometry is evaluated inside the integral. Only first reference
basis derivatives and first derivatives of the map are needed. A constant
has zero right-hand side, hence zero projected derivative. For affine maps,
physical differentiation stays in the supported space and this projection
recovers the existing exact polynomial derivative matrix.

In the one-dimensional example \cref{eq:curved-derivative-example}, the
projected first derivative $z_h$ is characterized by
\begin{equation}
 \int_0^1\widehat z_h\widehat\phi_a
       (1+\varepsilon-2\varepsilon\xi)\,d\xi
 =\int_0^1\widehat\phi_a\,d\xi.
\end{equation}
The Jacobian cancels on the right, but remains in the mass matrix on the left.
This illustrates why a physical projection is different from an unweighted
reference projection. To make the result explicit, set $\varepsilon=1/2$
and use the reference linear basis $(1,\xi)$. Then
\begin{equation}
 \begin{pmatrix}1&5/12\\5/12&1/4\end{pmatrix}
 \begin{pmatrix}Z_0\\Z_1\end{pmatrix}
 =\begin{pmatrix}1\\1/2\end{pmatrix},
 \qquad
 \widehat z_h(\xi)=\frac6{11}+\frac{12}{11}\xi.
 \label{eq:curved-projection-worked-example}
\end{equation}
This line approximates $1/(3/2-\xi)$ in the physical integral norm.
At $\xi=1$ it has value $18/11$, whereas the exact derivative is $2$.
The mismatch is compatible with the best-approximation property: endpoint
agreement was never imposed. This intentionally coarse example illustrates
the definition, rather than predicting the error of a refined simulation.

In actual assembly the integrals are replaced by quadrature, giving the
projection associated with that discrete inner product. Increasing
quadrature accuracy reduces integration error, but cannot make an
unrepresentable derivative belong to the fixed solution space.

\subsection{Recursive higher derivatives: the precise approximation}
\label{sec:curved-recursive-derivatives}
After the first operation, only the coefficients of $z_x=
\Pi_K^k\partial_xu_h$ are retained as the derivative state. A second
$x$-operation differentiates this represented approximation and projects
again:
\begin{equation}
 z_x=\Pi_K^k\partial_xu_h,\qquad
 z_{xx}=\Pi_K^k\partial_x z_x
       =\Pi_K^k\partial_x(\Pi_K^k\partial_xu_h).
\end{equation}
It does not return to $u_h$ and evaluate its exact second physical
derivative. The distinction persists even if both projections are assembled
with exact integration. The derivative graph implements this sequence by
reusing the same first-derivative matrices:
\begin{equation}
 U^{(0)}=U,\qquad U^{(\alpha)}=D_{K,j}U^{(\alpha-e_j)},
 \label{eq:curved-recursion}
\end{equation}
For $\alpha\ne0$, the code chooses $j=\max\{s:\alpha_s>0\}$.
The coordinate unit multi-index $e_j$ has a one in position $j$ and zeros
elsewhere, so $\alpha-e_j$ denotes the already computed parent state with
one fewer derivative. $U^{(\alpha)}$ is the coefficient vector of
$\mathscr D_K^\alpha u_h$, not a power of $U$. For $(1,1)$, for example,
the parent is $(1,0)$ and the last operation is $D_{K,y}$; for $(2,1)$,
the chain is $(0,0)\to(1,0)\to(2,0)\to(2,1)$. This gives exactly the
ordering in \cref{eq:projected-multi-index-order}.
For instance, the second operation in directions $i$ then $j$ represents
$\Pi_K^k\partial_{x_j}(\Pi_K^k\partial_{x_i}u_h)$. It generally differs
from projecting the exact second derivative once. Writing $\Pi=\Pi_K^k$,
\begin{equation}
 \Pi\partial_{x_j}(\Pi\partial_{x_i}u_h)
 -\Pi\partial_{x_j}\partial_{x_i}u_h
 =\Pi\partial_{x_j}(\Pi-I)\partial_{x_i}u_h.
 \label{eq:curved-recursion-defect}
\end{equation}
To read this identity, set $e_i=(\Pi-I)\partial_{x_i}u_h$.
The first stored derivative is $\partial_{x_i}u_h+e_i$, so the next
operation is $\Pi\partial_{x_j}\partial_{x_i}u_h+
\Pi\partial_{x_j}e_i$. The right-hand side of
\cref{eq:curved-recursion-defect} is precisely this last term.
A small error in the first projection does not by itself guarantee a small
higher-derivative error, because differentiating the error can amplify it.
There are also two distinct accuracy questions: error relative to derivatives
of the represented $u_h$, and error relative to derivatives of an underlying
smooth exact solution $u$. The latter includes differentiation of the input
approximation error $u_h-u$ as well as the intermediate projections.
It explains the difference between successive projection and a single
projection of the exact derivative. To isolate the ordering issue, subtract
the same identity with $i$ and $j$ interchanged. Exact mixed partials cancel:
\begin{align}
 (\mathscr D_{K,j}\mathscr D_{K,i}
   -\mathscr D_{K,i}\mathscr D_{K,j})u_h
 &=\Pi\partial_{x_j}(\Pi-I)\partial_{x_i}u_h\notag\\
 &\quad-\Pi\partial_{x_i}(\Pi-I)\partial_{x_j}u_h.
 \label{eq:curved-commutator}
\end{align}
Thus any failure to commute comes from the intermediate projection errors,
not from a failure of the physical chain rule on curved geometry. The two
terms vanish when the first derivatives are represented exactly in the
space. On curved elements they need not vanish, although particular states
or maps can still give commuting results.

The derivative graph stores one state per multi-index. On the affine path,
any valid parent gives the same exact derivative; the curved extension
retains one fixed parent as a definition of the approximation. It neither
computes every ordering nor verifies that alternative parents agree.
A multi-index remains a valid storage label with this convention, but it is
not evidence of order independence. Averaging derivative permutations or
projecting exact higher derivatives would define different sensor operators
and would require separate implementation and validation. Second- and
third-derivative convergence must therefore be assessed using
manufactured solutions. Projection alone is not a proof of their convergence.

The two routes are contrasted in \cref{fig:projected-derivative-route}.
\begin{figure}[htbp]
\centering
\begin{tikzpicture}[>=Latex,font=\small,
 state/.style={draw,rounded corners,fill=blue!4,minimum height=10mm,
               minimum width=21mm,inner sep=4pt}]
 \node[state] (u) at (0,1.65) {$u_h$};
 \node[state] (first) at (3.7,1.65) {$\partial_{x_i}u_h$};
 \node[state] (second) at (8,1.65) {$\partial_{x_j}\partial_{x_i}u_h$};
 \draw[->,thick] (u)--node[above] {$\partial_{x_i}$}(first);
 \draw[->,black!55,dashed,thick] (first)--node[above] {$\partial_{x_j}$}(second);
 \node[state,fill=teal!6] (projected) at (3.7,-.25) {$\Pi\partial_{x_i}u_h$};
 \node[state,fill=teal!6] (recursive) at (8,-.25)
    {$\Pi\partial_{x_j}(\Pi\partial_{x_i}u_h)$};
 \draw[->,teal!80!black,thick] (first)--node[left] {project: $\Pi$}(projected);
 \draw[->,teal!80!black,thick] (projected)--node[above]
    {$\Pi\partial_{x_j}$}(recursive);
 \node[align=left,anchor=east] at (1.9,-.25)
    {implemented route\\back into $V^k(K)$};
 \node[align=center] at (8,-1.55)
    {generally $\neq\Pi\partial_{x_j}\partial_{x_i}u_h$};
\end{tikzpicture}
\caption{Exact differentiation and the implemented recursive construction on a
curved element, with $\Pi=\Pi_K^k$. The first physical derivative generally
leaves $V^k(K)$. Projecting it back changes the input to the next derivative.
The lower route therefore differs from differentiating twice and projecting
only at the end. Each lower state is represented in the same finite element
space; no second derivative of the geometry map is needed by this route.}
\label{fig:projected-derivative-route}
\end{figure}


\subsection{Curved faces, normal directions, and length scales}
For a reference face with outward unit normal $\widehat{\boldsymbol n}$, the
physical surface factor and outward normal are
\begin{equation}
 s_F=\|w_KJ_K^{-\mathsf T}\widehat{\boldsymbol n}\|,
 \qquad \boldsymbol n_K=
 \frac{w_KJ_K^{-\mathsf T}\widehat{\boldsymbol n}}{s_F},
 \qquad ds=s_F\,d\widehat s.
 \label{eq:curved-face-metric}
\end{equation}
These quantities vary along a curved face. For face quadrature weights
$\omega_q$, use
\begin{equation}
 |F|_Q=\sum_q\omega_qs_{F,q},\qquad
 \widetilde\omega_q=\frac{\omega_qs_{F,q}}{|F|_Q},\qquad
 \sum_q\widetilde\omega_q=1.
 \label{eq:curved-face-normalization}
\end{equation}
Hence the implemented derivative-jump average is
\begin{equation}
 J_{F,c}^{(l)}\approx\sum_{|\alpha|=l}\sum_q\widetilde\omega_q
 \left| (\mathscr D_{K^-}^{\alpha}u_{h,c}^-)(x_q)
       -(\mathscr D_{K^+}^{\alpha}u_{h,c}^+)(x_q)\right|.
\end{equation}
Both traces are evaluated at the same physical point, using their respective
element maps. A constant jump has the same normalized magnitude on a straight
or curved face. The one-sided wave radii use $\boldsymbol n_K(x_q)$ at these
points, rather than a single centre normal.

Volume quadrature gives $|K|_Q=\sum_q\omega_qw_K(\xi_q)$ and the OFDG length
$h_K=(|K|_Q/|\widehat K|)^{1/d}$. The global mean and amplitude calculation
uses the same physical volume interpretation. KXRCF uses its separately defined
sampled radius and restricts its signed face integral to quadrature points
with incoming transport. Its denominator uses the corresponding physical
inflow measure. Thus sensing, normalization, and propagation directions all
refer to the curved physical geometry.

The algebra above uses exact integrals. The implementation assembles them
with positive, overintegrated quadrature. Consistent quadrature on the nested
spaces gives the analogous discrete projection identities, up to linear-solve
roundoff; it preserves the quadrature-defined physical means. Agreement with
exact integrals, derivative accuracy, and geometric sampling invariance still
require quadrature and refinement checks. This solution-space
projection does not convert rational NURBS geometry to polynomial geometry;
that separate, explicit geometry approximation retains the support boundary
stated in \cref{sec:support-boundary}.
